% ==============================================================================
% Abstract
% ==============================================================================

\begin{abstract}
Reading comprehension and vocabulary acquisition remain persistent challenges in higher education, particularly when students encounter domain-specific terminology in academic texts. While augmented reality (AR), intelligent tutoring systems (ITS), and gamification have each demonstrated positive effects on learning outcomes individually, no existing mobile platform integrates all three approaches into a unified, contextually-aware reading companion. This paper presents AR-DeC (Augmented Reality Didactic Companion), a novel mobile application that combines on-device optical character recognition (OCR), adaptive intelligent tutoring, AR-based visual overlays, and gamification mechanics to support vocabulary acquisition during natural reading activities. AR-DeC enables students to scan printed or digital text using their smartphone camera, whereupon the system automatically identifies unfamiliar or domain-specific words, generates contextual explanations adapted to the learner's current knowledge state via Bayesian Knowledge Tracing, and renders interactive AR annotations directly over the source material. A gamification engine incorporating points, badges, streaks, and adaptive quizzes sustains engagement and motivation through the ARCS (Attention, Relevance, Confidence, Satisfaction) motivational framework. The system architecture leverages React Native with Expo for cross-platform deployment, Google ML Kit for on-device OCR processing, ViroReact for AR rendering, and Firebase for backend services. We describe the theoretical foundations, system design, and technical implementation of AR-DeC, and outline a Design-Based Research (DBR) evaluation methodology involving 30--50 university students over a four-week intervention period. Planned assessments include pre-/post-vocabulary tests, the System Usability Scale (SUS), ARCS motivation questionnaires, and interaction log analysis. This work contributes a theoretically-grounded, technically-detailed blueprint for integrating AR, ITS, and gamification in mobile reading support tools, addressing a significant gap in the educational technology literature.
\end{abstract}
